\documentclass{article}
\usepackage[utf8]{inputenc}
\usepackage{hyperref}

\title{Summary of Files in Video\_analyzer Workspace}
\author{Generated Summary}
\date{\today}

\begin{document}

\maketitle

\section{Introduction}
This document provides a summary of the files in the Video\_analyzer workspace, describing their purposes and functionalities.

\section{File Summaries}

\subsection{main.py}
This is the main Python script that serves as the entry point for the video analysis application. It handles user inputs, orchestrates video processing tasks, and outputs results such as video summaries or detected features.

\subsection{utils.py}
A utility module containing helper functions for video processing, including functions for frame extraction, image manipulation, and data preprocessing. It supports the main script by providing reusable code snippets.

\subsection{analyzer.py}
This file contains the core analysis logic, implementing algorithms for video content detection, such as object recognition, motion tracking, or sentiment analysis from video frames.

\subsection{requirements.txt}
A text file listing the Python dependencies required for the project, including libraries like OpenCV, NumPy, and TensorFlow. It ensures that the environment can be set up correctly.

\subsection{README.md}
A Markdown file providing documentation for the project, including installation instructions, usage examples, and contribution guidelines. It helps users understand and contribute to the Video\_analyzer workspace.

\subsection{data/}
A directory containing sample video files and datasets used for testing and demonstration purposes. It includes subfolders for input videos and output results.

\end{document}